\section{Requirements of an algorithm implementation workflow}
\label{sec:requirements}

The design of \pqlang{} starts from a researcher implementing an
oracle-based scientific-computing algorithm. The researcher may begin at
the access-model, transform, or solver layer, retaining abstract
dependencies until their realizations are selected. This section derives
the implementation requirements of that workflow.  We first describe the canonical shape of a quantum
scientific-computing program (\secref{sec:qsc-structure}) and then extract
eight requirements (\secref{sec:qsc-requirements}), each of which will
reappear as a concrete language mechanism in \secref{sec:language} through
\secref{sec:ode}.

\subsection{The shape of a quantum scientific-computing program}
\label{sec:qsc-structure}

Across the literature on quantum linear algebra and differential equations,
programs have a remarkably stable layered shape, illustrated in
\figref{fig:stack}:

\begin{enumerate}
\item \textbf{Mathematical problem.}  A linear system $Ax=b$; an initial-value
  problem $u' = G u$, $u(0)=u_0$; a nonlinear polynomial system $u' =
  \sum_p F_p\, u^{\otimes p}$; or a PDE before spatial discretization.
\item \textbf{Input model (oracle paradigm).}  The problem is accessed through
  a declared access model: a block-encoding of $A$ or $G$ with scale factor
  $\alpha$~\cite{gilyen2019qsvt}; sparse-access oracles that locate and
  return nonzero entries; a QRAM serving an XOR database
  $|a\rangle|d\rangle \mapsto |a\rangle|d \oplus M[a]\rangle$
  \cite{giovannetti2008qram}; or a state-preparation oracle for $b$ or
  $u_0$.  Algorithm analyses are always relative to this model, and the
  model's parameters ($\alpha$, sparsity, width) enter the complexity bound
  \cite{dervovic2018qlss}.
\item \textbf{Primitive algebra.}  From the input model, larger objects are
  built: linear combinations of unitaries (LCU) with PREPARE--SELECT
  pairs~\cite{berry2015taylor,low2019qubitization}, tensor products,
  Kronecker sums, dilations of rectangular or non-Hermitian matrices,
  projectors, and shifted lattices arising from finite-difference stencils.
\item \textbf{Quantum transforms.}  Qubitization and QSVT
  \cite{low2019qubitization,gilyen2019qsvt} consume block-encodings to apply
  polynomial transformations of the encoded operator: matrix inversion
  filters for QLSS, Jacobi--Anger branches for Hamiltonian simulation,
  eigenstate filters, fixed-point search reflections.  Alternatively,
  product formulas or truncated-series LCUs realize the same functions with
  different cost/precision trade-offs.
\item \textbf{Solver assembly.}  Transforms are wrapped into solver
  skeletons: phase estimation readouts, amplitude-amplification boosts, or,
  in the ODE/PDE setting, solver families that reduce evolution to
  Hamiltonian-simulation calls---LCHS and its weighted quadrature of
  Hermitian branches~\cite{an2026lchs}, Schr\"odingerization's Fourier-lift
  embedding~\cite{jin2024schrodingerization}, contour decompositions,
  Carleman lifts of nonlinear systems into block-triangular linear
  ones~\cite{liu2021pnas}, or implicit-Euler history systems solved by a QLSS
  in one shot~\cite{berry2014ode,costa2022discrete}.
\item \textbf{Readout and norm recovery.}  Finally, since the solution is a
  quantum state, the program ends with a \emph{host-side} plan: phase or
  amplitude estimation, sample-based norm recovery through an independent
  matrix-norm probe, classical post-processing such as continued fractions,
  and statistics over repetitions.  Measurement is not part of the unitary
  algorithm.
\end{enumerate}

\begin{figure}[t]
  \centering
  \begin{tikzpicture}[
      font=\small,
      layer/.style={draw=black!60, rounded corners=2pt, fill=black!3,
        text width=10.9cm, minimum height=1.05cm, align=center},
      side/.style={font=\footnotesize\itshape, text=black!60, align=left},
      arr/.style={-{Stealth[length=2.2mm]}, black!60},
    ]
    \node[layer, fill=black!8] (prob) at (0,0)
      {\textbf{Mathematical problem}\\[-1pt]{\footnotesize $Ax=b$;\; $u'=Gu$;\; $u'=\sum_p F_p u^{\otimes p}$;\; PDEs before discretization}};
    \node[layer] (input) at (0,-1.55)
      {\textbf{Input model (oracle paradigm)}\\[-1pt]{\footnotesize block-encoding($\alpha$);\; sparse location/entry;\; QRAM XOR database;\; state preparation}};
    \node[layer] (prim) at (0,-3.1)
      {\textbf{Primitive algebra}\\[-1pt]{\footnotesize LCU / PREPARE--SELECT;\; tensor \& Kronecker sums;\; dilation;\; projectors;\; stencil shifts}};
    \node[layer] (trans) at (0,-4.65)
      {\textbf{Quantum transforms}\\[-1pt]{\footnotesize qubitization \& QSVT;\; QSP phase design;\; product formulas;\; truncated Taylor series}};
    \node[layer] (solve) at (0,-6.2)
      {\textbf{Solver assembly}\\[-1pt]{\footnotesize QLSS;\; LCHS / Schr\"odingerization / contour;\; Carleman lift;\; history systems;\; amplitude amplification}};
    \node[layer, fill=black!8] (read) at (0,-7.75)
      {\textbf{Host-side readout \& recovery}\\[-1pt]{\footnotesize phase/amplitude estimation;\; norm probes;\; post-processing;\; statistics}};

    \draw[arr] (prob) -- (input);
    \draw[arr] (input) -- (prim);
    \draw[arr] (prim) -- (trans);
    \draw[arr] (trans) -- (solve);
    \draw[arr] (solve) -- (read);

    \node[side, anchor=west] at (5.95, -1.55) {R1, R2};
    \node[side, anchor=west] at (5.95, -3.10) {R3, R5};
    \node[side, anchor=west] at (5.95, -4.65) {R4, R6};
    \node[side, anchor=west] at (5.95, -6.20) {R6, R7, R8};
  \end{tikzpicture}
  \caption{The canonical layering of a quantum scientific-computing program.
  Each layer is accessed through the layer below it by a well-defined
  interface; the right margin annotates which of the requirements
  R1--R8 of \secref{sec:qsc-requirements} the layer stresses.  The top and
  bottom layers (shaded) are host-side concerns; the language must connect to
  them without embedding them in the quantum IR.}
  \label{fig:stack}
\end{figure}

Two features of this shape drive everything that follows.  First, the
interfaces \emph{between} layers are mathematical contracts, and the objects
that realize a given contract are interchangeable: the same block-encoding
slot may be filled by a gate network, a QRAM lookup, an alias-sampling
construction, or an application-specific reversible circuit, with different
resource profiles but identical semantics at the contract level.  Second, the
shape is \emph{deep and parameterized}: a PDE discretization parameter $n$
enters as register widths and repetition counts, not as a longer gate list,
and the scientifically meaningful object is the symbolic composition, not any
particular expansion.

\subsection{Requirements}
\label{sec:qsc-requirements}

From this structure we extract requirements.  Throughout,
Table~\ref{tab:requirements} summarizes the requirement, its mathematical
source, the deficiency it exposes in common practice, and the mechanism that
\pqlang{} builds for it.

\paragraph{R1: Deferred oracle realization.}
Algorithms are published relative to abstract oracles; implementations are
chosen later, and the choice is a deployment decision (gate synthesis
vs.\ hardware QRAM vs.\ native operator), not an algorithmic one.  A language
must therefore let a program be \emph{complete as an algorithm} while its
oracle slots are empty---validatable in that state, shareable in that state,
and closable without rewriting.  Python operation factories already defer construction in many SDKs.
The additional requirement here is that a partially implemented program
survives storage and restoration without the original factory.

\paragraph{R2: Realization independence.}
For a given oracle slot, several concrete realizations with different
resource/cost profiles must be interchangeable behind one semantic contract,
including QRAM-backed ones (where the data table itself is a runtime
resource, not part of the program text).  This is what makes an algorithm
library reusable across input regimes and what makes cost studies---switching
one slot from QROM to alias sampling, say---into a one-line change.

\paragraph{R3: Compositional structure preservation.}
Subroutine calls, static repetition, control, and adjoint blocks must survive
end-to-end: in the IR, in its serialization, and in exports.  Eager flattening
destroys the object to which precision annotations, resource accounts, and
review attention attach; it also makes program size explode with scale
parameters long before execution is ever attempted.  The IR, not the
frontend, is the artifact of record.

\paragraph{R4: Adjoint and control as tracked capabilities.}
Scientific-computing programs use $U^{\dagger}$ and controlled-$U$ pervasively
(uncompute-and-prepare patterns, oblivious amplitude amplification,
eigenstate filtering).  Whether a composed object admits adjoint or control
is a compositional property of its parts; a language must infer and check it
\emph{at composition time}, before any backend is asked to realize it.

\paragraph{R5: Symbolic structure.}
Repetition counts, addressing widths, and truncation degrees are program
parameters, not reasons to replicate text.  The frontend must never expand
them; cost models, not materialization, must guard execution; and even
textual export must remain logarithmic in repetition counts.  Scale
invariants must be part of the specification, not folklore:
``serialization must not unconditionally replicate a repeated body'' is a
testable statement.

\paragraph{R6: Promise and precision provenance.}
Block-encoding factors $\alpha$, spectral-norm promises
($\|A\|\le a$, $\sigma_{\min}(A)\ge\sigma$), dissipativity of
the ODE generator, degree choices of polynomial approximations, and known
approximation caveats (finite quadrature, omitted contour poles, Carleman
tails) are all part of a program's meaning.  None of them is a gate; none of
them is safely left in a paper's appendix.  A language must carry them as
first-class metadata attached to the objects they qualify, with explicit
provenance, and must not pretend they are guarantees: a promise declared
without evidence is recorded as a declaration.  Conversely, the language
core need not, and should not, reify a precision target $\varepsilon$ as a
type: precision is produced by generator functions and applications, and the
language's job is to keep its provenance auditable.

\paragraph{R7: Readout separation.}
The unitary algorithm and the measurement plan are different objects with
different lifetimes.  In this framework, separating measurement plans from coherent operations
makes solver outputs reusable by later quantum subroutines.  Measurement, reset, post-selection, and statistics belong
to a host-side plan that references, but does not contaminate, the IR.

\paragraph{R8: Library extensibility without core changes.}
The algorithm zoo grows by research, not by committee: LCHS, Schr\"odingerization,
and Carleman-based pipelines were research artifacts within recent memory.
A language in which adding a solver family means modifying the grammar, the
IR, the serializer, or a central type registry will always lag the
literature.  New algorithms must be expressible as new libraries---new
protocols, new input models, new generator functions---over a fixed core.

\begin{table}[t]
  \centering
  \caption{Requirements derived from the structure of quantum scientific
  computing, common-practice deficiencies, and the corresponding \pqlang{}
  mechanisms.  Each mechanism is detailed in the referenced section.}
  \label{tab:requirements}
  \small
  \setlength{\tabcolsep}{4pt}\begin{tabular}{@{}p{0.45cm}p{3.0cm}p{3.4cm}p{6.2cm}@{}}
    \toprule
    \tabhead{ } & \tabhead{Requirement} & \tabhead{Implementation concern} & \tabhead{\pqlang{} mechanism (section)} \\
    \midrule
    R1 & Deferred oracle realization & Preserving an incomplete algorithm beyond its Python generator & Open modules: declared access roles with absent bodies; validation and serialization work while open (\secref{sec:oracles}) \\
    R2 & Realization independence & Realization choice hard-wired into algorithm code & One oracle slot, many bindings: gate networks, QRAM-typed resources, alias sampling; partial and repeated binding via the \texttt{bind} linking pass (\secref{sec:oracles}) \\
    R3 & Structure preservation & Frontends flatten to gate lists; exports lose call structure & Module-preserving RIR: \texttt{Call}/\texttt{Repeat}/\texttt{Control}/\texttt{Adjoint} stay symbolic in IR, JSON, and export (\secref{sec:rir}) \\
    R4 & Tracked adjoint/control & Adjoint/control availability discovered (or not) at backend & Compositional capability inference checked against contextual usage at validation time (\secref{sec:capabilities}) \\
    R5 & Symbolic structure & Counts materialized eagerly; programs explode with system size & \texttt{Repeat} nodes preserved with static counts; logarithmic export; cost-model guards; unrolling forbidden in generation and serialization (\secref{sec:rir}) \\
    R6 & Promise provenance & $\alpha$, norm bounds, dissipativity, approximation caveats live in papers, not code & First-class attributes ($\alpha$, \texttt{SpectralPromise} with evidence strings, dissipativity flags, per-method pending items) attached to modules and contracts (\secref{sec:oracles}, \secref{sec:ode}) \\
    R7 & Readout separation & Maintaining coherent solver outputs for later subroutines & Readout plans separate from quantum operations; restricted memory effects remain explicit (\secref{sec:readout}) \\
    R8 & Extensibility & Adding algorithms and input roles over existing primitives & Library-level structural interfaces and algorithm contracts; new roles do not require a global registry (\secref{sec:protocols}) \\
    \bottomrule
  \end{tabular}
\end{table}

\paragraph{Resulting programming paradigm.}
These requirements fix a three-phase programming paradigm that recurs
throughout \pqlang{}:
\begin{quote}
\emph{Generate} programs in Python against declared input models;
\emph{store and share} them as RIR---open or bound;
\emph{export or execute} on backends, with readout planned at the host.
\end{quote}
Within generation, the differential-equation stack
(\secref{sec:ode}) adds a sharper discipline that we call
\emph{problem--adapter--protocol}: the \emph{problem layer} keeps the
mathematical input paradigm native to the application (a PDE, a polynomial
ODE, a sparse matrix); \emph{adapters} explicitly construct the objects an
algorithm needs (block-encodings, lifted initial states); and \emph{protocols}
are Python functions that generate new oracles---whose own implementations
may still be open.  RIR stores the result: generated modules plus the
not-yet-implemented oracle slots they depend on.  The division of
responsibility is explicit: the language is responsible for composability;
precision is the responsibility of generator functions and applications
(R6).
